%! Describing the Distribution of a Quantitative Variable — Unit 1, Lesson 1.4 — Student Packet Cover Sheet
\documentclass[12pt]{article}
\usepackage{apstats-article}
\usepackage{apstats-boxes}
\usepackage{ltablex}
\keepXColumns

\begin{document}

% Full-bleed navy banner. \coverbanner MEASURES the title block and sizes the
% band to it, so a lesson title long enough to wrap grows the band rather than
% running out the bottom of it.
\coverbanner{Unit 1}{Lesson 1.4 \quad Describing the Distribution of a Quantitative Variable}

\namedateperiod
\vspace{0.18in}

\boxguard[14]
\begin{learningtargetbox}
\small
% GRADUAL RELEASE: the vocabulary is taught in the guided notes, so the targets name it
% plainly. The AP Learning Objective (UNC-1.H) stays in the teacher-facing plan.
\begin{enumerate}[leftmargin=1.5em, itemsep=5pt]
  \item \ldots{} explain why the \textbf{center} of a distribution, reported by itself, is not a
        description of anything.
  \item \ldots{} name the \textbf{shape} of a distribution --- \textbf{symmetric},
        \textbf{skewed left}, or \textbf{skewed right} --- and say why the \emph{tail} decides
        that name, not the pile.
  \item \ldots{} point out the \textbf{unusual features} of a display --- \textbf{outliers},
        \textbf{gaps}, and \textbf{clusters} --- and say why an outlier is investigated, not deleted.
  \item \ldots{} describe a distribution completely --- \textbf{shape, unusual features,
        center, variability} --- all four in context, with the variable and its units named.
\end{enumerate}
\end{learningtargetbox}

\vspace{0.14in}

\boxguard[14]
\begin{tocbox}
\small
\renewcommand{\arraystretch}{1.5}
\begin{tabularx}{\linewidth}{@{} c l X r @{}}
\rowcolor{sky}
\textbf{\#} & \textbf{Component} & \textbf{Description} & \textbf{Score} \\
\midrule
1 & Warm-Up & Two pictures of one data set, a middle value, and one dot sitting far from
                 the rest & \blank{1.2cm} \\
2 & Guided Notes \& Practice & Two checkout lines with the same middle wait, three leaning
                 pictures, and the four things every description says & \blank{1.2cm} \\
% AP Practice is assigned but NOT scored: its score cell prints NA, never a \blank{}.
% Row 4 is the lesson's GRADED practice. Paper homework is the DEFAULT for this course; on the
% occasions DeltaMath is assigned instead, that replaces row 4 and its score cell is struck.
3 & AP Practice & Commute times for 40 employees: AP-style multiple choice and one
                 free-response set --- practice, not a grade & \textbf{NA} \\
4 & Homework & Match, correct a classmate, describe one distribution in full, and one
                 multiple-choice item --- \textbf{scored, due the first class after two study
                 halls} & \blank{1.2cm} \\
\midrule
  & \multicolumn{2}{r}{\textbf{Total}} & \blank{1.2cm} \\
\end{tabularx}
\end{tocbox}

\vspace{0.14in}

\boxguard[8]
\begin{remindbox}
\small
Two checkout lines in these notes have \emph{exactly the same} middle wait, and one of them you
would never stand in again --- a center by itself is a summary of nothing. A distribution takes
\textbf{four} things to describe --- shape, unusual features, center, and spread --- each in
context, with the variable and its units named. \textbf{A skew is named for the side its tail
runs out on}, not the side the data piles up on. And a value far from the rest is a finding to
investigate, never an error to delete.
\end{remindbox}

\end{document}
